% ថ្នាក់ទី ១១ (វិទ្យាសាស្ត្រ) · ជំពូកទី ១ · មេរៀនទី ៦ · សៀវភៅទំព័រ ៤៤–៥៦
\chapter{វិចារអនុមានរួមគណិតវិទ្យា}

\begin{objBox}[វត្ថុបំណង]
\begin{enumerate}[1]
  \item កំណត់វិចារអនុមានរួមគណិតវិទ្យា
  \item ការប្រើអនុមានរួមគណិតវិទ្យា
  \item ទ្រឹស្តីបទទ្វេធា ។
\end{enumerate}
\end{objBox}

\section{គោលការណ៍នៃវិចារអនុមានរួមគណិតវិទ្យា}

គេគ្រាន់តែសង្កេតមើលឧទាហរណ៍ខ្លះៗនៃផលបូកចំនួនគត់សេស
\[
\begin{aligned}
&\text{ផលបូក } 1 \text{ តួតាងដោយ } P(1) = 1 = 1^2\\
&\text{ផលបូក } 2 \text{ តួតាងដោយ } P(2) = 1 + 3 = 2^2\\
&\text{ផលបូក } 3 \text{ តួតាងដោយ } P(3) = 1 + 3 + 5 = 3^2
\end{aligned}
\]
គេអាចទាញបានផលបូក $n$ តួតាងដោយ $P(n) = 1 + 3 + 5 + \ldots + (2n - 1) = n^2$ ។

បន្ទាប់ពីបានរូបមន្តនេះមក គេត្រូវបញ្ជាក់ថារូបមន្ត $P(n)$ ពិតចំពោះគ្រប់ចំនួនគត់ $n$ ។ ក្នុងករណីនេះគេផ្តើមផ្ទៀងផ្ទាត់ $P(n)$ ដោយឱ្យ $n = 1$ ។

បើ $n = 1$ \quad $1 = 1^2$ ដូច្នេះ $P(1)$ ពិត ។

ឧបមាថា $P(n)$ ពិត ចំពោះតម្លៃ $n$ ណាមួយ ។ គេត្រូវស្រាយថា $P(n)$ ពិតនាំឱ្យបាន $P(n+1)$ ក៏ពិតដែរ ។

បើ $P(n)$ ពិត គេបាន $1 + 3 + 5 + \ldots + (2n - 1) = n^2$ ។ ដើម្បីស្រាយបញ្ជាក់ថា $P(n+1)$ ពិតលុះត្រាតែគេបាន $P(n+1) = (n+1)^2$ ។

\textbf{របៀបស្រាយបញ្ជាក់}

គេដឹងថា
\[
\begin{aligned}
P(n+1) &= 1 + 3 + 5 + \ldots + (2n - 1)\\
&\quad + [2(n + 1) - 1]\\
&= n^2 + (2n + 2 - 1)\\
&= n^2 + 2n + 1 = (n + 1)^2
\end{aligned}
\]
$P(n+1)$ ពិត ។ ដូច្នេះ $P(n)$ ពិតចំពោះគ្រប់ $n \in \mathbb{N}$ ។

\begin{theorem}[គោលការណ៍នៃវិចារអនុមានរួមគណិតវិទ្យា]
វិចារអនុមានរួម $P(n)$ ជាសំណើដែលទាក់ទងនឹងចំនួនគត់ $n$ ។ ដើម្បីស្រាយបញ្ជាក់ថា $P(n)$ ពិតចំពោះគ្រប់ $n \in \mathbb{N}$ គេត្រូវ
\begin{enumerate}[1]
  \item ផ្ទៀងផ្ទាត់ថា $P(n)$ ពិតចំពោះ $n = 1$
  \item ឧបមាថា $P(n)$ ពិតចំពោះតម្លៃ $n$
  \item ស្រាយបញ្ជាក់ថា $P(n)$ ពិតនាំឱ្យបាន $P(n+1)$ ពិត ។
\end{enumerate}
\end{theorem}

\begin{example}[លំហាត់គំរូ ១]
ស្រាយបញ្ជាក់សមភាព $1 \times 2 + 2 \times 3 + 3 \times 4 + \ldots + n \times (n + 1) = \dfrac{1}{3}n(n + 1)(n + 2)$ ពិតចំពោះគ្រប់ $n \in \mathbb{N}$ ដោយប្រើវិចារអនុមានរួមគណិតវិទ្យា ។
\end{example}
\ans
ចំពោះ $n = 1$ ៖ អង្គខាងឆ្វេង $1 \times 2 = 2$ \quad អង្គខាងស្តាំ $\dfrac{1}{3} \times 1 \times (1 + 1)(1 + 2) = 2$ ។

នោះអង្គខាងឆ្វេង $=$ អង្គខាងស្តាំ ។ ដូច្នេះ ចំពោះ $n = 1$ សមភាពពិត ។

ឧបមាថា សមភាពពិតចំពោះ $n = k$ គេបាន
\[
\begin{gathered}
1 \times 2 + 2 \times 3 + 3 \times 4 + \ldots + k \times (k + 1) = \frac{1}{3}k(k + 1)(k + 2)\\
(k = 1, 2, 3, \ldots) \text{ ។}
\end{gathered}
\]
បង្ហាញថាសមភាពពិតចំពោះ $n = k$ នាំឱ្យបានសមភាពពិតចំពោះ $n = k + 1$
\[
\begin{aligned}
&1 \times 2 + 2 \times 3 + 3 \times 4 + \ldots + k(k + 1) + (k + 1)(k + 2)\\
&= \frac{1}{3}(k + 1)(k + 2)(k + 3)
\end{aligned}
\]
គេមាន $1 \times 2 + 2 \times 3 + 3 \times 4 + \ldots + k \times (k + 1) = \dfrac{1}{3}k(k + 1)(k + 2)$ ។ ថែម $(k + 1)(k + 2)$ លើអង្គទាំងពីរនៃសមីការ គេបាន
\[
\begin{aligned}
&1 \times 2 + 2 \times 3 + 3 \times 4 + \ldots + k \times (k + 1) + (k + 1)(k + 2)\\
&= \frac{1}{3}k(k + 1)(k + 2) + (k + 1)(k + 2)\\
&= (k + 1)(k + 2)\left(\frac{1}{3}k + 1\right)\\
&= (k + 1)(k + 2)\left(\frac{k + 3}{3}\right)\\
&= \frac{1}{3}(k + 1)(k + 2)(k + 3)
\end{aligned}
\]
ពិត ។ ដូច្នេះ សមភាពពិតចំពោះគ្រប់ $n \in \mathbb{N}$ ។

ដូច្នេះ $1 \times 2 + 2 \times 3 + 3 \times 4 + \ldots + n \times (n + 1) = \dfrac{1}{3}n(n + 1)(n + 2)$ \quad $(n = 1, 2, 3, \ldots)$ ។

\begin{example}[លំហាត់គំរូ ២]
ចំពោះ $h > 0$ និងចំនួនគត់ធម្មជាតិ $n \ge 2$ បង្ហាញថា $(1 + h)^n > 1 + nh$ ដោយប្រើវិចារអនុមានរួមគណិតវិទ្យា ។
\end{example}
\ans
ចំពោះ $n = 2$ គេបាន $(1 + h)^2 = 1 + 2h + h^2 > 1 + 2h$ ព្រោះ $h^2 > 0$ ។ ដូច្នេះ វិសមភាពនេះពិត ចំពោះ $n = 2$ ។

ឧបមាថាវិសមភាពនេះពិត ចំពោះ $n = k$ $(k \ge 2)$ នោះ $(1 + h)^k > 1 + kh$ ពិត ។ គេនឹងបង្ហាញថាវាពិតដល់ $n = k + 1$ $(k \ge 2)$ គឺ
\[ (1 + h)^{k+1} > 1 + (k + 1) \times h \quad (1) \text{ ។} \]
គេមាន $(1 + h)^k > 1 + kh$ ដោយ $h > 0$ នោះ $1 + h > 0$ ។ គុណអង្គទាំងពីរនៃ $(1)$ និង $1 + h > 0$ គេបាន $(1 + h)^k(1 + h) > (1 + kh)(1 + h)$
\[
\begin{aligned}
(1 + h)^{k+1} &> 1 + kh + h + kh^2\\
&= 1 + (k + 1)h + kh^2 > 1 + (k + 1)h
\end{aligned}
\]
ព្រោះ $kh^2 > 0$ ។ ដូច្នេះ វិសមភាព $(1 + h)^n > 1 + nh$ ពិតចំពោះគ្រប់ $(n \ge 2)$ ។

\begin{example}[លំហាត់គំរូ ៣]
ស្រាយបញ្ជាក់ថា $n^3 + 2n$ ជាពហុគុណនៃ $3$ ចំពោះគ្រប់ចំនួនគត់ $n$ ដោយប្រើវិចារអនុមានរួមគណិតវិទ្យា ។
\end{example}
\ans
ចំពោះ $n = 0$ គេបាន $0$ ជាពហុគុណនៃ $3$ ពិត ។

ឧបមាថាសមភាពពិតចំពោះចំនួនគត់ $n = k$ គេបាន $k^3 + 2k$ ជាពហុគុណនៃ $3$ ពិត យើងបាន $k^3 + 2k = 3q$, $q$ ជាចំនួនគត់ ។

គេបង្ហាញថាវាពិតចំពោះ $n = k + 1$ គឺ $(k + 1)^3 + 2(k + 1)$ ជាពហុគុណនៃ $3$ ។ គេបាន
\[
\begin{aligned}
(k + 1)^3 + 2(k + 1) &= k^3 + 3k^2 + 3k + 1 + 2k + 2\\
&= k^3 + 3k^2 + 5k + 3\\
&= (k^3 + 2k) + 3 \times (k^2 + k + 1)\\
&= 3q + 3(k^2 + k + 1) = 3(q + k^2 + k + 1)
\end{aligned}
\]
មានន័យថា $(k + 1)^3 + 2(k + 1)$ ជាពហុគុណនៃ $3$ ។

ដូច្នេះ $n^3 + 2n$ ជាពហុគុណនៃ $3$ ចំពោះគ្រប់ចំនួនគត់ $n$ ។

\begin{pro}
\begin{enumerate}[a]
  \item ស្រាយបញ្ជាក់វិសមភាព $\left(1 + \dfrac{1}{2} + \dfrac{1}{3} + \ldots + \dfrac{1}{n}\right)(1 + 2 + 3 + \ldots + n) \ge n^2$ ដោយប្រើវិចារអនុមានរួមគណិតវិទ្យា ។
  \item ស្រាយបញ្ជាក់ថា $2^{n+1} + 3^{2n-1}$ ជាពហុគុណនៃ $7$ ចំពោះគ្រប់ $n \in \mathbb{N}$ ដោយប្រើវិចារអនុមានរួមគណិតវិទ្យា ។
  \item ស្រាយបញ្ជាក់សមភាព $1^2 + 2^2 + 3^2 + \ldots + n^2 = \dfrac{n(n + 1)(2n + 1)}{6}$ ដោយប្រើវិចារអនុមានរួមគណិតវិទ្យា ។
\end{enumerate}
\end{pro}

\section{ទ្រឹស្តីបទទ្វេធា}

ពេលគេពន្លាតស្វ័យគុណនៃទ្វេធា គេបានពហុធា ឬស៊េរីមួយ ។ វាមិនមែនជាស៊ីតធរណីមាត្រ ឬស៊ីតនព្វន្តឡើយ តែវាមានលំនាំគំរូមួយច្បាស់លាស់ ។ ក្នុងមេរៀននេះគេនឹងសិក្សារូបមន្តពន្លាតនៃស្វ័យគុណ នៃទ្វេធាឱ្យបានឆាប់រហ័ស ។ ជាដំបូងពិនិត្យមើលពន្លាតនៃកន្សោម $(x + y)^n$ ចំពោះតម្លៃខ្លះនៃ $n$ ។

\begin{example}[ឧទាហរណ៍ ១]
\[
\begin{aligned}
(x + y)^0 &= 1\\
(x + y)^1 &= 1x + 1y\\
(x + y)^2 &= 1x^2 + 2xy + 1y^2\\
(x + y)^3 &= 1x^3 + 3x^2y^1 + 3x^1y^2 + 1y^3\\
(x + y)^4 &= 1x^4 + 4x^3y^1 + 6x^2y^2 + 4x^1y^3 + 1y^4\\
(x + y)^5 &= 1x^5 + 5x^4y^1 + 10x^3y^2\\
&\quad + 10x^2y^3 + 5x^1y^4 + 1y^5\\
&\ \,\vdots
\end{aligned}
\]
ក្នុងជួរដេកនីមួយៗកើតមានលំនាំគំរូដូចៗគ្នា ។
\end{example}

ពីឧទាហរណ៍ខាងលើ គេសង្កេតក្នុងជួរដេកនីមួយៗនៃពន្លាតកន្សោម $(x + y)^n$ ឃើញថា ៖
\begin{itemize}
  \item ពន្លាតទ្វេធានីមួយៗជាពហុធាដែលមាន $(n + 1)$ តួ ។
  \item ក្នុងពន្លាតទ្វេធានីមួយៗស្វ័យគុណនៃ $x$ ចុះ $1$ ដឺក្រេជាបន្តបន្ទាប់ និងស្វ័យគុណនៃ $y$ កើន $1$ ដឺក្រេជាបន្តបន្ទាប់បើគេអានពីឆ្វេងទៅស្តាំ ។
  \item ផលបូកនៃស្វ័យគុណក្នុងតួនីមួយៗនៃពហុធាស្មើ $n$ ។
\end{itemize}
\begin{itemize}
  \item លេខមេគុណចាប់ផ្តើមនិងចុងក្រោយគឺ $1$ ដូចគ្នា ហើយផលបូកនៃលេខមេគុណពីរតួជាប់គ្នានៃពន្លាត $(x + y)^3$ ជាលេខមេគុណនៃពន្លាត $(x + y)^4 \ldots$ ។
\end{itemize}
\[
\begin{array}{ccccccccc l}
 & 1 & & 3 & & 3 & & 1 & & \quad\text{លេខមេគុណនៃ } (x+y)^3\\
1 & & 4 & & 6 & & 4 & & 1 & \quad\text{លេខមេគុណនៃ } (x+y)^4
\end{array}
\]
ឬលេខមេគុណនៃតួផ្សេងទៀតកំណត់បានអាស្រ័យនឹងលេខមេគុណនៃតួបន្ទាប់ខាងមុខគឺ
\[
\begin{gathered}
\text{លេខមេគុណនៃតួទី } k + 1\\
= \frac{(\text{លេខមេគុណនៃតួទី } k) \times (\text{ស្វ័យគុណនៃ } x \text{ ក្នុងតួទី } k)}{k} \text{ ។}
\end{gathered}
\]

\begin{example}[ឧទាហរណ៍ ២]
\[
\begin{aligned}
(x + y)^5 = 1x^5 + 5x^4y^1 + 10x^3y^2 &\\
&\quad + 10x^2y^3 + 5xy^4 + 1y^5
\end{aligned}
\]
ដែល $5 = \dfrac{1 \times 5}{1}$, \quad $10 = \dfrac{5 \times 4}{2}$, \quad $10 = \dfrac{10 \times 3}{3}$, \quad $5 = \dfrac{10 \times 2}{4}$, \quad $1 = \dfrac{5 \times 1}{5}$ ។
\end{example}

ពីឧទាហរណ៍ ១ បើគេផ្តាច់យកតែលេខមេគុណនិងតម្រៀបលេខទាំងនេះជារាងត្រីកោណ គេឃើញមានលំនាំគំរូជាច្រើន ។ គេបាន
\[
\begin{array}{ccccccccccc}
 & & & & & 1 & & & & & \\
 & & & & 1 & & 1 & & & & \\
 & & & 1 & & 2 & & 1 & & & \\
 & & 1 & & 3 & & 3 & & 1 & & \\
 & 1 & & 4 & & 6 & & 4 & & 1 & \\
1 & & 5 & & 10 & & 10 & & 5 & & 1
\end{array}
\]
តម្រៀបលេខមេគុណទាំងនេះជារាងត្រីកោណហៅថា \textbf{ត្រីកោណប៉ាស្កាល់} និងមានសារៈសំខាន់ណាស់សម្រាប់រកលេខមេគុណក្នុងពន្លាតទ្វេធា $(x + y)^n$ ដែលមានស្វ័យគុណតូចតែវាពិបាកអនុវត្តខ្លាំងណាស់ចំពោះការពន្លាតទ្វេធា $(x + y)^n$ ដែលមានស្វ័យគុណធំ ។

ចំណែកឯស៊ីត $1, 1, 2, 3, 5, 8, 13, 21, \ldots$ (ផលបូកតាមអង្កត់ទ្រេតនៃត្រីកោណប៉ាស្កាល់) ហៅថា \textbf{ស៊ីត Fibonacci} ។

\begin{example}[លំហាត់គំរូ ៤]
ប្រើត្រីកោណប៉ាស្កាល់ពន្លាតកន្សោម $(x + y)^5$ ។
\end{example}
\ans
ដោយប្រើជួរដេកទី $5$ នៃត្រីកោណប៉ាស្កាល់ គេបានលេខមេគុណ
\[
\begin{array}{ccccccccccc}
 & 1 & & 4 & & 6 & & 4 & & 1 & \\
1 & & 5 & & 10 & & 10 & & 5 & & 1
\end{array}
\]
គេបាន $(x + y)^5 = 1x^5 + 5x^4y^1 + 10x^3y^2 + 10x^2y^3 + 5xy^4 + 1y^5$ ។

\begin{pro}
ប្រើត្រីកោណប៉ាស្កាល់ពន្លាតកន្សោម $(x + y)^6$ ។
\ans
$(x + y)^6 = x^6 + 6x^5y + 15x^4y^2 + 20x^3y^3 + 15x^2y^4 + 6xy^5 + y^6$ ។
\end{pro}

គេចាំបាច់កំណត់គំរូតួទូទៅនៃពន្លាត $(x + y)^n$ គឺ
\[
\begin{aligned}
(x + y)^n = 1x^n + ?x^{n-1}y^1 + ?x^{n-2}y^2 + \ldots &\\
&\quad + ?x^{n-r}y^r + \ldots + 1y^n
\end{aligned}
\]
\[
\begin{array}{c|c}
\text{តួ} & \text{លេខមេគុណ}\\ \hline
1 & 1\\[4pt]
2 & \dfrac{1 \times n}{1}\\[8pt]
3 & \dfrac{n(n - 1)}{1 \times 2}\\[6pt]
\vdots & \vdots\\
(r + 1) & \dfrac{n(n - 1)(n - 2)\ldots(n - r + 1)}{1 \times 2 \times 3 \times \ldots \times r}\\[6pt]
\vdots & \vdots\\
(n + 1) & 1
\end{array}
\]
ដូច្នេះ
\[
\begin{aligned}
(x + y)^n &= 1x^n + \frac{n}{1}x^{n-1}y^1 + \frac{n(n - 1)}{1 \times 2}x^{n-2}y^2 + \ldots\\
&\quad + \frac{n(n - 1)(n - 2)\ldots(n - r + 1)}{1 \times 2 \times 3 \times \ldots \times r}x^{n-r}y^r + \ldots + 1y^n \text{ ។}
\end{aligned}
\]

ពីត្រីកោណប៉ាស្កាល់ គេឃើញថាលេខមេគុណនៃទ្វេធាជាចំនួនដែលរៀបឆ្លុះគ្នា ហើយគេដឹងថា
\[
\begin{aligned}
1 &= \frac{n!}{n!\,0!} = {}_nC_0 = {}_nC_n\\
\frac{n}{1} &= \frac{n!}{(n - 1)!\,1!} = {}_nC_1 = {}_nC_{n-1}\\
\frac{n(n - 1)}{1 \times 2} &= \frac{n!}{(n - 2)!\,2!} = {}_nC_2 = {}_nC_{n-2}\\
\frac{n(n - 1)(n - 2)}{1 \times 2 \times 3} &= \frac{n!}{(n - 3)!\,3!} = {}_nC_3 = {}_nC_{n-3}\\
&\ \,\vdots\\
\frac{n(n - 1)(n - 2)\ldots(n - r + 1)}{1 \times 2 \times 3 \times \ldots \times r} &= \frac{n!}{(n - r)!\,r!} = {}_nC_r = {}_nC_{n-r}
\end{aligned}
\]
ចំនួនថេរ ${}_nC_0, {}_nC_1, {}_nC_2, \ldots, {}_nC_n$ ហៅថាលេខមេគុណនៃទ្វេធា $(x + y)^n$ ។ គេបាន
\[
\begin{aligned}
(x + y)^n &= {}_nC_0x^n + {}_nC_1x^{n-1}y^1\\
&\quad + {}_nC_2x^{n-2}y^2 + \ldots\\
&\quad + {}_nC_rx^{n-r}y^r + \ldots + {}_nC_ny^n
\end{aligned}
\]
ហៅថា \textbf{ទ្រឹស្តីបទទ្វេធា} ។

\begin{theorem}[ទ្រឹស្តីបទទ្វេធា]
\[
\begin{aligned}
(x + y)^n &= \sum_{r=0}^{n} {}_nC_r\,x^{n-r}y^r\\
&= {}_nC_0\,x^n + {}_nC_1\,x^{n-1}y^1\\
&\quad + {}_nC_2\,x^{n-2}y^2 + \ldots + {}_nC_n\,y^n \text{ ។}
\end{aligned}
\]
\end{theorem}

\begin{example}[លំហាត់គំរូ ៥]
ដោយប្រើជួរដេកទី $6$ នៃត្រីកោណប៉ាស្កាល់គណនាលេខមេគុណនៃទ្វេធា ${}_6C_0$, ${}_6C_1$, ${}_6C_2$, ${}_6C_3$, ${}_6C_4$, ${}_6C_5$, ${}_6C_6$ ។
\end{example}
\ans
\[
\begin{array}{ccccccccccccc}
 & 1 & & 5 & & 10 & & 10 & & 5 & & 1 & \\
1 & & 6 & & 15 & & 20 & & 15 & & 6 & & 1\\
{}_6C_0 & & {}_6C_1 & & {}_6C_2 & & {}_6C_3 & & {}_6C_4 & & {}_6C_5 & & {}_6C_6
\end{array}
\]
ដូច្នេះ ${}_6C_0 = 1$, ${}_6C_1 = 6$, ${}_6C_2 = 15$, ${}_6C_3 = 20$, ${}_6C_4 = 15$, ${}_6C_5 = 6$, ${}_6C_6 = 1$ ។

\begin{pro}
\begin{enumerate}[1]
  \item ដោយប្រើទ្រឹស្តីបទទ្វេធាចូរពន្លាតកន្សោមខាងក្រោម ៖
  \begin{tasks}(4)
    \task $(1 + 3x)^4$
    \task $(2x - y)^5$
    \task $\left(a + \dfrac{1}{a^2}\right)^6$
    \task $(1 + x)^n$ ។
  \end{tasks}
  \item កំណត់តួទី $4$ នៃពន្លាតទ្វេធា $(x + y)^{13}$ ។
\end{enumerate}
\ans
\begin{enumerate}[1]
  \item
  \begin{enumerate}[a]
    \item $(1 + 3x)^4 = 1 + 12x + 54x^2 + 108x^3 + 81x^4$
    \item $(2x - y)^5 = 32x^5 - 80x^4y + 80x^3y^2 - 40x^2y^3 + 10xy^4 - y^5$
    \item $\left(a + \dfrac{1}{a^2}\right)^6 = a^6 + 6a^3 + 15 + \dfrac{20}{a^3} + \dfrac{15}{a^6} + \dfrac{6}{a^9} + \dfrac{1}{a^{12}}$
    \item $(1 + x)^n = 1 + {}_nC_1\,x + {}_nC_2\,x^2 + {}_nC_3\,x^3 + \ldots + {}_nC_n\,x^n$ ។
  \end{enumerate}
  \item ${}_{13}C_3\,x^{10}y^3 = 286x^{10}y^3$ ។
\end{enumerate}
\end{pro}

\section*{មេរៀនសង្ខេប ជំពូកទី ១}
\begin{note}[មេរៀនសង្ខេប]
\begin{enumerate}[1]
  \item ស៊ីតនៃចំនួនពិត គឺជាអនុគមន៍លេខដែលកំណត់ពី $\mathbb{N}$ ទៅ $\mathbb{R}$ ។
  \item ស៊ីតនព្វន្ត គឺជាស៊ីតនៃចំនួនពិតដែលមានតួនីមួយៗ (ក្រៅពីតួទី $1$) ស្មើនឹងតួមុនបន្ទាប់បូកចំនួនថេរ $d$ មួយ (ហៅថាផលសងរួម $d = a_{n+1} - a_n$) ។
  \item រូបមន្តតួទូទៅនៃស៊ីតនព្វន្ត $a_n = a_1 + (n - 1) \times d$ ។
  \item ផលបូក $n$ តួនៃស៊ីតនព្វន្ត $S_n = \dfrac{(a_1 + a_n) \times n}{2}$ ។
  \item ស៊ីតធរណីមាត្រ គឺជាស៊ីតចំនួនពិតដែលតួនីមួយ (ក្រៅពីតួទី $1$) ស្មើនឹងតួមុនបន្ទាប់គុណនឹងចំនួនថេរ $q$ ដែល $q \ne 0$ ។ ចំនួនថេរ $q$ នេះហៅថាផលធៀបរួមស៊ីតធរណីមាត្រ ។
  \item រូបមន្តតួទូទៅនៃស៊ីតធរណីមាត្រ $a_{n+1} = a_n \times q$, \ $a_n = a_1 \times q^{n-1}$ ។
  \item ផលបូក $n$ តួដំបូងនៃស៊ីតធរណីមាត្រ $S_n = \dfrac{a_1(q^n - 1)}{q - 1}$ ដែល $q \ne 1$ ។
  \item ផលបូកស៊ីតធរណីមាត្រអនន្តតួ $S_\infty = \dfrac{a_1}{1 - q}$ បើ $|q| < 1$ ។
  \item រូបមន្តផលបូកស៊ីតមានលំដាប់ខ្ពស់ \
  $\displaystyle\sum_{k=1}^{n} k = \frac{1}{2}n(n + 1)$, \ $\displaystyle\sum_{k=1}^{n} k^2 = \frac{1}{6}n(n + 1)(2n + 1)$, \ $\displaystyle\sum_{k=1}^{n} k^3 = \left[\frac{1}{2}n(n + 1)\right]^2$ ។
  \item ស៊ីត $(b_n)$ ជាផលសងតួលំដាប់ទី $1$ នៃស៊ីត $(a_n)$ ដែល $a_n = a_1 + \displaystyle\sum_{k=1}^{n-1} b_k$, \ $n \ge 2$ ។
  \item ទំនាក់ទំនងកំណើនទម្រង់ $a_{n+1} = a_n + pn$ នោះផលសងស៊ីត $(b_n)$ ដែល $b_n = a_{n+1} - a_n = pn$ និង $a_n = a_1 + \displaystyle\sum_{k=1}^{n-1} b_k = a_1 + \sum_{k=1}^{n-1} pk$ ។
  \item ចំពោះ $n \ge 2$ គេបាន $S_n - S_{n-1} = a_n$ និង $S_1 = a_1$ ។
  \item ទំនាក់ទំនងកំណើននៃស៊ីត $(a_n)$ ទម្រង់ $a_{n+2} + pa_{n+1} + qa_n = 0$ មានសមីការសម្គាល់ $x^2 + px + q = 0$ ដែលឫស $\alpha$ និង $\beta$ បំពេញលក្ខខណ្ឌ $\alpha + \beta = -p$, $\alpha\beta = q$ នោះគេបាន $a_{n+2} - (\alpha + \beta)a_{n+1} + \alpha\beta a_n = 0$ អាចបំប្លែងជា $a_{n+2} - \alpha a_{n+1} = \beta(a_{n+1} - \alpha a_n)$, \ $a_{n+2} - \beta a_{n+1} = \alpha(a_{n+1} - \beta a_n)$ ។ ទំនាក់ទំនងទាំង $2$ នេះគឺជាស៊ីតធរណីមាត្រ ។
  \item សំណើ $P(n)$ ពិតចំពោះគ្រប់ចំនួនគត់ធម្មជាតិ $n$ ហើយសំណើរខាងក្រោម $(i)$ និង $(ii)$ ពិតគឺ \\
  $(i)$ \ $P(n)$ ពិតចំពោះ $n = 1$ \\
  $(ii)$ \ ឧបមាថា $P(n)$ ពិតចំពោះគ្រប់ចំនួនគត់ធម្មជាតិ $n = k$ រួចបង្ហាញថា $P(n+1)$ ពិតចំពោះគ្រប់ $n = k + 1$ ។
  \item ទ្រឹស្តីបទទ្វេធា $(x + y)^n = \displaystyle\sum_{r=0}^{n} {}_nC_r\,x^{n-r}y^r = {}_nC_0\,x^n + {}_nC_1\,x^{n-1}y^1 + {}_nC_2\,x^{n-2}y^2 + \ldots + {}_nC_n\,y^n$ ។
\end{enumerate}
\end{note}

\section*{លំហាត់ជំពូកទី ១}
\begin{enumerate}[1]
  \item តួទី $5$ នៃស៊ីតនព្វន្តស្មើ $14$ និងតួទី $12$ ស្មើ $35$ ។ គណនាផលសងរួម តួទី $1$ និងតួទី $n$ ។
  \item ផលបូកនៃ $n$ តួដំបូងនៃស៊ីតមួយកំណត់ដោយ $S_n = n^2 + n$ ។ គណនាតួទី $5$ នៃស៊ីត និងបង្ហាញថាតួនៃស៊ីតនេះជាស៊ីតនព្វន្ត ។
  \item តួទី $3$ នៃស៊ីតធរណីមាត្រស្មើ $3$ និងតួទី $6$ ស្មើនឹង $\dfrac{3}{8}$ ។ គណនាផលធៀបរួមតួទី $1$ និងផលបូក $8$ តួដំបូង ។
  \item គណនា $a$ ដើម្បីឱ្យបីចំនួន $(a - 3)$, $(a + 1)$ និង $(4a - 2)$ ជាបីតួតគ្នានៃស៊ីតធរណីមាត្រ ។
  \item គេឱ្យបួនចំនួនជាស៊ីតធរណីមាត្រ ។ គណនាតួទី $2$ និងទី $3$ នៃស៊ីត បើគេដឹងថាតួទី $1$ ស្មើនឹង $6$ ហើយតួទី $4$ ស្មើនឹង $162$ ។
  \item គេទិញផ្ទះមួយតម្លៃ $30000$ ដុល្លារ ។ អ្នកទិញត្រូវបង់ប្រាក់ជា $6$ ដំណាក់កាល ដែលបង្កើតបានជាស៊ីតធរណីមាត្រមានផលធៀបរួម $\dfrac{9}{10}$ ។ កំណត់ប្រាក់ដែលត្រូវបង់តាមដំណាក់កាលនីមួយៗ ។
  \item គេមានស៊ីតនព្វន្ត $(U_n)$ និង $(V_n)$ កំណត់ដោយ $U_1 = 3$ និងផលសងរួមស្មើ $3$ និង $V_1 = 4$ និងផលសងរួមស្មើ $2$ ។
  \begin{enumerate}[a]
    \item កំណត់ $U_n$ និង $V_n$ ជាអនុគមន៍នៃ $n$ ។
    \item គេមានស៊ីត $(W_n)$ កំណត់ដោយ $W_n = \dfrac{U_n}{V_n}$ ។ បង្ហាញថាស៊ីត $(W_n)$ ជាស៊ីតកើន និងកំណត់ $n$ ដើម្បីឱ្យ $W_n - W_{n-1} = \dfrac{1}{28}$ ។
  \end{enumerate}
  \item កំណត់ផលសងរួមនិងតួទី $1$ $U_1$ នៃស៊ីតនព្វន្ត $(U_n)_{n \in \mathbb{N}}$ កើននិងដឹងថា $\begin{cases} U_1^2 + U_2^2 + U_3^2 = 98\\ U_1 - U_2 + U_3 = 4 \end{cases}$ ។ គណនា $U_{10}$ និងផលបូក $10$ តួដំបូងនៃស៊ីត ។
  \item តួដំបូងនៃស៊ីតធរណីមាត្រស្មើ $3$ ហើយមានផលធៀបរួម $r$ ដែលតម្លៃដាច់ខាតនៃ $r$ តូចជាង $1$ ។ ផលបូក $3$ តួដំបូងស្មើ $\dfrac{8}{9}$ នៃផលបូក $6$ តួដំបូង ។ គណនាផលបូកស៊ីតធរណីមាត្រអនន្តតួ ។
  \item
  \begin{enumerate}[1]
    \item ដោះស្រាយក្នុង $\mathbb{R}$ សមីការ $3x^2 - 8x + 4 = 0$ \ $(E)$ ។
    \item គេមាន $(U_n)_{n \in \mathbb{N}}$ ជាស៊ីតធរណីមាត្រកើនដែលមាន $U_3$ និង $U_4$ នៃស៊ីតនេះជាចម្លើយនៃសមីការ $(E)$ ។
    \begin{tasks}(2)
      \task ទាញរកផលធៀបរួមនៃស៊ីតនេះ ។
      \task គណនាតួទី $1$ $U_1$ នៃស៊ីត ។
    \end{tasks}
  \end{enumerate}
  \item សរសេរបីតួដំបូងនៃស៊ីតនព្វន្តដោយដឹងថា $S_{10} = 210$ និង $S_{20} = 820$ ។
  \item គេដឹងថាផលបូកតួទី $1$ និងតួទី $4$ នៃស៊ីតនព្វន្តស្មើនឹង $2$ និងផលបូកការេរបស់វាស្មើនឹង $20$ ។ គណនាផលបូកប្រាំបីតួដំបូងនៃស៊ីត ។
  \item គេមាន $S_m$ និង $S_n$ ជាផលបូក $m$ តួដំបូង និង $n$ តួដំបូងរៀងគ្នានៃស៊ីតនព្វន្តមួយដែល $\dfrac{S_m}{S_n} = \dfrac{m^2}{n^2}$; $(n \ne m)$ ។ តាងតួទី $m$ គឺ $a_m$ និងតួទី $n$ គឺ $a_n$ ។ បង្ហាញថា $\dfrac{a_m}{a_n} = \dfrac{2m - 1}{2n - 1}$ ។
  \item គេមានស៊ីតធរណីមាត្រ $12, 4, \dfrac{4}{3}, \ldots$ ។
  \begin{enumerate}[a]
    \item គណនាតួទី $10$
    \item តើចំនួន $\dfrac{4}{729}$ ជាតួទីប៉ុន្មាននៃស៊ីត?
    \item គណនាផលបូក $20$ តួដំបូងនៃស៊ីតធរណីមាត្រ ។
  \end{enumerate}
  \item គេឱ្យ $(a_n)$ ជាស៊ីតធរណីមាត្រ បើគេដឹងថា $a_n = 2(3)^{n-1}$ ។ គណនា $S_n$ ។
  \item គណនាផលបូក $1 + 2x + 3x^2 + \ldots + (n - 1)x^{n-2} + nx^{n-1}$ ។
  \item គណនាតួទី $1$ នៃស៊ីតធរណីមាត្រអនន្តតួដែលមាន $q = \dfrac{3}{5}$ និង $S_\infty = 40$ ។
  \item គេឱ្យបីចំនួនជាស៊ីតធរណីមាត្រ ។ គណនាចំនួនទាំងនេះបើគេដឹងថាផលគុណនៃចំនួនទាំងនោះស្មើនឹង $3375$ ហើយផលបូកវាស្មើនឹង $93$ ។
  \item គណនាផលបូក $n$ តួដំបូងនៃស៊ីតនីមួយៗខាងក្រោម
  \begin{enumerate}[a]
    \item ស៊ីត $(a_n)$ ៖ $1, \dfrac{1}{1 + 2}, \dfrac{1}{1 + 2 + 3}, \ldots, \dfrac{1}{1 + 2 + 3 + 4 + \ldots + n}$
    \item ស៊ីត $(b_n)$ ៖ $\dfrac{1}{(1 \times 3)^2}, \dfrac{2}{(3 \times 5)^2}, \dfrac{3}{(5 \times 7)^2}, \ldots, \dfrac{n}{[(2n - 1)(2n + 1)]^2}$ ។
  \end{enumerate}
  \item
  \begin{enumerate}[a]
    \item កំណត់តួទី $n$ នៃស៊ីត $1, 2, 6, 15, 31, 56, \ldots$
    \item គណនាផលបូក $n$ តួដំបូងនៃស៊ីតនេះ ។
  \end{enumerate}
  \item
  \begin{enumerate}[a]
    \item គណនា $\displaystyle\sum_{k=1}^{n} (2k^2 - 1)$
    \item ដោយប្រើសំណួរ ក. គណនាផលបូក $1 + 7 + 17 + 31 + \ldots + 799$ ។
  \end{enumerate}
  \item សរសេរផលបូក $1 + 4 + 7 + 10 + 13 + \ldots + 298$ ដោយប្រើ $\Sigma$ ។
  \item ដោយប្រើវិចារអនុមានរួមគណិតវិទ្យាស្រាយបញ្ជាក់ថា $\displaystyle\sum_{k=1}^{n} 2^{k-1} = 2^n - 1$ ។
  \item គេមានស៊ីត $(U_n)$ ដែល $n$ ជាចំនួនគត់កំណត់ដោយ $U_{n+1} = 2U_n + 1$ និង $U_0 = 1$ ហើយស៊ីត $(V_n)$ កំណត់ដោយ $V_n = U_n + 1$ ។
  \begin{enumerate}[a]
    \item បង្ហាញថាស៊ីត $(V_n)$ ជាស៊ីតធរណីមាត្រ ។
    \item ទាញរក $U_n$ ជាអនុគមន៍នៃ $n$ ។
    \item សិក្សាភាពម៉ូណូតូននៃស៊ីត $(U_n)$ ។
    \item ចំពោះគ្រប់ចំនួនគត់ $n$ គណនាផលបូក $S_n = U_0 + U_1 + U_2 + \ldots + U_n$ ។
  \end{enumerate}
  \item គេមាន $(a_n)$ កំណត់ដោយទំនាក់ទំនងកំណើន $a_1 = 1$, $a_{n+1} = \dfrac{1}{2}a_n + n$ ។
  \begin{enumerate}[a]
    \item តាង $b_n = 2^na_n$ ។ កំណត់តួទី $n$ នៃស៊ីត $(b_n)$ ។
    \item កំណត់តួទី $n$ នៃស៊ីត $(a_n)$ ។
  \end{enumerate}
  \item
  \begin{enumerate}[a]
    \item សរសេរ $(x^2 - 2y)^7$ ដោយប្រើ $\Sigma$
    \item កំណត់តួទី $6$ នៃពន្លាត $(x^2 - 2y)^7$ ។
  \end{enumerate}
  \item បង្ហាញថាផលបូកនៃលេខមេគុណក្នុងពន្លាត $(1 + x)^n$ គឺ $2^n$ ។
  \item បង្ហាញថា ${}_nC_0 + {}_nC_2 + {}_nC_4 + \ldots + {}_nC_{n-1} = {}_nC_1 + {}_nC_3 + {}_nC_5 + \ldots + {}_nC_n = 2^{n-1}$ ដែល ${}_nC_0, {}_nC_1, {}_nC_2, \ldots, {}_nC_n$ ជាលេខមេគុណក្នុងពន្លាត $(1 + x)^n$ ដែល $n$ ជាចំនួនគត់សេស ។
  \item គេឱ្យស៊ីត $(a_n)$ មួយកំណត់ដោយ $a_1 = 1$, $a_{n+1} = \dfrac{a_n}{\sqrt{1 + 168(a_n)^2}}$, $n \in \mathbb{N}$ ។
  \begin{enumerate}[a]
    \item តាង $b_n = \dfrac{1}{(a_n)^2}$ ។ បង្ហាញថា $(b_n)$ ជាស៊ីតនព្វន្តមួយ
    \item គណនា $b_n$, $S_n = \displaystyle\sum_{k=1}^{n} b_k$ និង $a_n$ ជាអនុគមន៍នៃ $n$ ។
  \end{enumerate}
  \item គេឱ្យពហុធា $f(x) = px^2 + qx$ ។
  \begin{enumerate}[a]
    \item កំណត់តម្លៃពិតនៃ $p$ និង $q$ ដើម្បីឱ្យ $f(x) = f(x - 1) + x$ គ្រប់ $x$
    \item ទាញរកផលបូក $\displaystyle\sum_{k=1}^{n} k = 1 + 2 + 3 + \ldots$ ។
  \end{enumerate}
  \item គេឱ្យស៊ីតនព្វន្តមួយ $a_1 = \dfrac{5\pi}{3}$ និងមានផលសងរួម $d = \dfrac{\pi}{6}$ ។
  \begin{enumerate}[a]
    \item គណនាតួទី $168$ នៃស៊ីតនេះ
    \item តើតួទីប៉ុន្មាននៃស៊ីតនេះស្មើនឹង $\dfrac{8\pi}{3}$ ។
  \end{enumerate}
  \item គេឱ្យសមីការ $x^2 + \alpha x + \beta$ ។ កំណត់តម្លៃ $\alpha$ និង $\beta$ ដើម្បីឱ្យផលបូកការេនៃឫសសមីការនេះស្មើនឹង $14$ ហើយ $1, \alpha, \beta$ ជាបីតួបន្តគ្នានៃស៊ីតនព្វន្តចុះជាច់ខាតមួយ ។
  \item គេឱ្យស៊ីត $(a_n)$ មួយកំណត់ដោយ $a_1 = 2$, $a_{n+1} = \dfrac{5a_n - 1}{a_n + 3}$ ។
  \begin{enumerate}[a]
    \item តាង $b_n = \dfrac{1}{a_n - 1}$ ។ បង្ហាញថា $(b_n)$ ជាស៊ីតនព្វន្តមួយ
    \item សរសេរ $b_n$ និង $a_n$ ជាអនុគមន៍នៃ $n$ ។
  \end{enumerate}
  \item គេឱ្យស៊ីត $(a_n)$ មួយកំណត់ដោយ $a_1 = 2$, $a_{n+1} = \sqrt[3]{2 + 3(a_n)^2}$, $n = 1, 2, 3, \ldots$
  \begin{enumerate}[a]
    \item ស្រាយថា $b_n = 1 + a_n^3$ ជាស៊ីតធរណីមាត្រ
    \item សរសេរ $b_n$ និង $a_n$ ជាអនុគមន៍នៃ $n$
    \item $S_n = b_1 + b_2 + b_3 + \ldots + b_n$ ។
  \end{enumerate}
  \item គណនាផលបូក $S_n = \dfrac{13}{10} + \dfrac{303}{100} + \dfrac{5003}{1000} + \ldots + \dfrac{(2n - 1)00\ldots003}{1000\ldots000}$ ។
  \item គេឱ្យ $(a_n)$ កំណត់ដោយ $\begin{cases} a_1 = a, \quad a_2 = b\\ a_n = \dfrac{a_{n-1} + a_{n-2}}{2} \end{cases}$ ។
  \begin{enumerate}[a]
    \item ស្រាយថា $b_n = a_n - a_{n-1}$ ជាស៊ីតធរណីមាត្រ
    \item សរសេរ $a_n$ ជាអនុគមន៍នៃ $a$, $b$ និង $n$ ។
  \end{enumerate}
  \item គេឱ្យស៊ីត $(a_n)$ មួយកំណត់លើ $\mathbb{N}$ ដោយ $\begin{cases} a_1 = 5 + 3\sqrt{5}\\ a_{n+1} = \dfrac{1}{2}\left(a_n + \dfrac{5}{a_n}\right) \end{cases}$ ។
  \begin{enumerate}[a]
    \item តាង $b_n = \dfrac{a_n - \sqrt{5}}{a_n + \sqrt{5}}$ ។ សរសេរទំនាក់ទំនងរវាង $b_{n+1}$ និង $b_n$
    \item សរសេរ $b_n$ និង $a_n$ ជាអនុគមន៍នៃ $n$
    \item បើ $c_n = 1 + b_n$ គណនា $P_n = c_1 \times c_2 \times c_3 \times \ldots \times c_n$ ។
  \end{enumerate}
  \item គេមានស៊ីតចំនួនពិត $(a_n)$ មួយកំណត់លើ $\mathbb{N}$ ដោយ $\begin{cases} a_1 = \sqrt{5}\\ a_{n+1} = a_n^2 - 2 \end{cases}$ ។
  \begin{enumerate}[a]
    \item បង្ហាញថា $a_n > 2$, $\forall n \in \mathbb{N}$
    \item តាង $b_n = a_n + \sqrt{a_n^2 - 4}$ សរសេរទំនាក់ទំនងរវាង $b_n$ និង $b_{n+1}$
    \item សរសេរ $a_n$ និង $b_n$ ជាអនុគមន៍ $n$ ។
  \end{enumerate}
  \item គេមានស៊ីត $(a_n)$ មួយកំណត់លើ $\mathbb{N}$ ដោយ $\begin{cases} a_1 = 1\\ a_{n+1} = \dfrac{1}{2}a_n - \dfrac{1}{3} \end{cases}$ ។
  \begin{enumerate}[a]
    \item បង្ហាញថាគេអាចកំណត់តម្លៃ $k$ ដើម្បីឱ្យ $b_n = 3a_n + k$ ជាស៊ីតធរណីមាត្រមួយគ្រប់ $\mathbb{N}$
    \item សរសេរ $b_n$ និង $a_n$ ជាអនុគមន៍នៃ $n$ ចំពោះតម្លៃ $k$ ខាងលើ ។
  \end{enumerate}
  \item គេឱ្យកន្សោម $S_n = \dfrac{1}{7} + \dfrac{2}{7^2} + \dfrac{3}{7^3} + \ldots + \dfrac{n - 1}{7^{n-1}} + \dfrac{n}{7^n}$ ។ គណនា $S_n$ ។
  \item គេមានស៊ីតចំនួនពិត $(a_n)$ មួយកំណត់លើ $\mathbb{N}$ ដោយ $\begin{cases} a_1 = \dfrac{1}{3}\\ a_{n+1} = \dfrac{1}{3}\left(1 + \dfrac{1}{n}\right)a_n \end{cases}$ ។
  \begin{enumerate}[a]
    \item សរសេរ $b_n$ ជាអនុគមន៍នៃ $n$
    \item តាង $S_n = a_1 + a_2 + a_3 + \ldots + a_n$ ។ គណនា $S_n - \dfrac{1}{3}S_n$ រួចទាញរក $S_n$ ជាអនុគមន៍នៃ $n$ ។
  \end{enumerate}
  \item កំណត់តួទី $n$, $a_n$ នៃស៊ីត $(a_n)$ ខាងក្រោម ៖
  \begin{tasks}(2)
    \task $\begin{cases} a_1 = 2\\ a_{n+1} = 2a_n + 2 \end{cases}$
    \task $\begin{cases} a_1 = \dfrac{13}{2}\\ a_{n+1} = \dfrac{1}{2}a_n + 3 \end{cases}$
    \task $\begin{cases} a_1 = 0\\ a_{n+1} = (\sqrt{3} + \sqrt{2})a_n + \sqrt{3} - \sqrt{2} + 1 \end{cases}$
  \end{tasks}
  \item កំណត់តួទី $n$, $a_n$ នៃស៊ីត $(a_n)$ ខាងក្រោម ៖
  \begin{tasks}(2)
    \task $\begin{cases} a_1 = 0\\ a_{n+2} = 2a_{n+1} + 3a_n \end{cases}$
    \task $\begin{cases} a_1 = 1\\ a_{n+2} = 6a_{n+1} - 9a_n \end{cases}$
    \task $\begin{cases} a_1 = 0, \quad a_2 = 1\\ a_{n+2} = \sqrt{3}a_{n+1} - 4a_n \end{cases}$ ។
  \end{tasks}
\end{enumerate}

\section*{លំហាត់អនុវត្តន៍}
\begin{enumerate}[1]
  \item ស្រាយបញ្ជាក់ថាចំពោះគ្រប់ $n \ge 1$, $1^2 + 2^2 + 3^2 + \ldots + n^2 = \dfrac{n(n + 1)(2n + 1)}{6}$ ដោយប្រើវិចារអនុមានរួមគណិតវិទ្យា ។
  \item បង្ហាញថាចំពោះគ្រប់ចំនួនគត់ $n$
  \begin{tasks}(2)
    \task ចំនួន $4^n + 2$ ចែកដាច់នឹង $3$ ។
    \task ចំនួន $3^{n+3} - 4^{4n+2}$ ចែកដាច់នឹង $11$ ។
  \end{tasks}
  \item គេមានស៊ីត $(a_n)$ ដែល $n$ ជាចំនួនគត់កំណត់ដោយ $a_{n+1} = \sqrt{a_n + 2}$ និង $a_0 = 0$ ។ បង្ហាញតាមវិចារអនុមានរួមគណិតវិទ្យាថា
  \begin{tasks}(2)
    \task ចំពោះគ្រប់ចំនួនគត់ $n$, $a_n \le 2$ ។
    \task ចំពោះគ្រប់ចំនួនគត់ $n$, $a_n \le a_{n+1}$ ។
  \end{tasks}
  \item បង្ហាញតាមវិចារអនុមានរួមគណិតវិទ្យាថា ចំពោះគ្រប់ចំនួនគត់ធម្មជាតិ $n$ \quad $(1 + x)^n \ge 1 + nx + \dfrac{n(n - 1)}{2}x^2$ ចំពោះគ្រប់ចំនួនពិត $x \ge 0$ ។
  \item ដោយប្រើទ្រឹស្តីបទទ្វេធា ចូរពន្លាតកន្សោមខាងក្រោម ៖
  \begin{tasks}(3)
    \task $(3x - 1)^4$
    \task $(2x + y)^6$
    \task $(a + b)^6$ ។
  \end{tasks}
  \item សរសេរពន្លាតកន្សោមខាងក្រោមដោយប្រើ $\Sigma$ ៖
  \begin{tasks}(3)
    \task $(x - y)^5$
    \task $(2x + y)^6$
    \task $(a + b)^{12}$ ។
  \end{tasks}
\end{enumerate}
